Figure~\ref{fig:equity_costs} provides a simple implementation-cost sensitivity.
Starting from the same out-of-sample target weights, we subtract 25 basis points
per unit of target-weight turnover. We then additionally charge a constant
30 basis points per year on short notional, a deliberately simple borrowing-cost
benchmark motivated by the order of magnitude of median shorting fees reported
by \citet{jensenkellymalamudpedersen2026}. The Sharpe ratios fall from
$3.20$, $3.35$, and $3.25$ to approximately $2.59$, $2.70$, and $2.61$
after commissions, and to $2.55$, $2.66$, and $2.57$ after both commissions
and the borrowing-cost adjustment. The borrowing-cost exercise is a sensitivity
analysis rather than a reconstruction of stock-specific historical lending fees.

\begin{figure}[!htbp]
    \centering
    \includegraphics[width=0.93\textwidth]{figures/equity_with_costs.pdf}
    \caption{Cumulative out-of-sample excess-return wealth before costs, after
    a 25 basis point commission per unit of target-weight turnover, and after
    additionally charging a constant 30 basis points per year on short
    notional. The borrowing fee is an illustrative sensitivity rather than a
    stock-specific historical lending-cost series.}
    \label{fig:equity_costs}
\end{figure}




